\section{Background} 
\subsection{GateKeeper Filtering Algorithm}

GateKeeper~\cite{Alser2017} is the first FPGA-based pre-alignment filter that significantly improves the filtering speed. The algorithm is built upon Shifted Hamming Distance~(SHD)~\cite{xin2015shifted} and mainly consists of bitwise AND, XOR, and shift operations. Its simple and bitwise nature makes it applicable for hardware acceleration.

GateKeeper starts by encoding the read and its candidate reference segment in 2 bits (i.e., A = 00, C = 01, G = 10, T = 11) to prepare the bitwise representations of the strings. Then, it performs the XOR operation between bitvectors of the read and the reference segment to prepare the Hamming mask for exact match detection. Differentiating characters that depict mismatches are indicated as \textit{`1'} on the resulting bitvector, and matching characters are shown with \textit{`0'}. The approximate matching phase begins if the error threshold~(i.e., \textit{e}) is more than 0. In a loop of incrementing the error tolerance invariant \textit{k = $\{1,2,...,e\}$} one by one at each iteration, intermediate bitvectors are prepared by shifting read bitvector to the right and left by \textit{k} bits for deletion and insertion, respectively.
At the end of one iteration, the shifted bitvector undergoes XOR operation with reference segment bitvector to detect possible errors, thus every iteration produces two masks: one for deletion and one for insertion. Since the XOR operation signifies the difference in 2-bits per encoded character comparison, every two-bit is combined with bitwise OR to simplify the differences on individual bitvectors and reduce resource usage. 

The final stage of approximate matching is an AND operation on all \textit{2e+1} masks. Nevertheless, a \textit{0}-bit at a particular position is dominant, and AND will yield a \textit{0}-bit indicating a match even if all of the other bitvectors have a \textit{1}-bit value at that position, which signals a mismatch.
To compensate for this issue, the bitvectors are \textit{amended} before AND to turn short streaks of \textit{0}s into \textit{1}s considering these \textit {0}s are useless and do not represent an informative part in Hamming masks. For amendment, after the AND operation, the errors are counted by following a window approach with a look-up table. The comparison pair is rejected if the number of \textit{1s} as errors in the final bitvector exceeds the error threshold and is accepted if otherwise. For a step-by-step example of the GateKeeper algorithm's workflow, please refer to Sup. Figure S.1.  

\subsection{Unified Memory Architecture}
GPU provides a massive amount of parallelism to accelerate the compute-intensive work, which otherwise takes a long time with the CPU. The specialized cores of GPU process the data that physically reside on the device, thus, utilizing GPU as an accelerator inevitably requires moving the data from CPU to GPU memory via PCI-e bus. Because this creates an extra task for GPU when compared to CPU-only execution, adopting an effective memory management and data movement strategy plays a critical role. CUDA framework offers \textit{unified memory} for creating a virtual space, which GPU and CPU have access with a single pointer~\cite{harris_2018_unified}.

Unified memory does not remove the necessity of data movement altogether but maintains an automatic migration of data on-demand, providing data locality~\cite{harris_2018_unified}. Conceptually, the unified memory model fits GateKeeper-GPU in many aspects. For instance, the reference genome's designated segments are requested only when the mapper signals a potential match for the specific read in seeding, which creates an on-demand access situation. Likewise, it is sufficient to copy a single read only once to GPU memory for its multiple candidate reference segments. In our experience with different memory allocation methods at the development stage of GateKeeper-GPU, we also achieved more efficient overall execution with unified memory than with pinned memory format. Therefore, we adopted a unified memory model in GateKeeper-GPU.

Along with unified memory abstraction, the CUDA framework introduces \textit{memory advise} and \textit{asynchronous memory prefetching} features for optimized data migration during execution. Data access patterns of an object, such as preferred location, can be specified by memory advice so that the processor favors the optimal placement of the data for migration decisions. In addition, asynchronous data prefetching initiates the data transfer to the device before the data is used for minimizing the overhead caused by page faults in runtime~\cite{chien2019performance}. Prefetching is supported by Pascal and later architectures with CUDA 8 in GPU.


\subsection{Related Work}
Finding optimal solutions for the approximate string matching problem has been the focus of genome sequence studies due to the nature of data and its tremendous size. One of the tools built for this purpose, Edlib~\cite{vsovsic2017edlib}, is a CPU framework that calculates edit distance~(also known as Levenstein Distance~\cite{levenshtein1966binary}) by utilizing Myers' bitvector algorithm~\cite{myers1999fast} for optimal pairwise sequence alignment. Since Edlib finds the exact edit distance, we hold Edlib's global alignment results as the ground truth for our accuracy analysis.

The verification stage of read mapping tools determines the exact similarity between the given two sequences as read and candidate reference segment, therefore, the efforts on pre-alignment filtering prioritize quickly eliminating pairs with an apparent dissimilarity. Shifted Hamming Distance~(SHD)\cite{xin2015shifted} filters out highly erroneous sequence pairs with a bit-parallel and SIMD-parallel algorithm. MAGNET~\cite{alser2017magnet} addresses some of SHD's shortcomings, such as not considering leading and trailing zeros in bitmasks and consecutive bit counting for edit calculation, which are the main sources of high false accept rate. With a new filtering strategy, it improves the filtering accuracy by up to two orders of magnitude. Shouji~\cite{alser2019shouji}, on the other hand, takes a different approach to pre-alignment filtering. It starts by preparing a neighborhood map for identifying the common substrings between two sequences within an error threshold. Then, it finds non-overlapping common subsequences with a sliding window approach and decides on accepting or rejecting the pair according to the length of common subsequences. With an FPGA design, Shouji can reduce sequence alignment duration up to $18.8\times$ with two orders of magnitude higher accuracy than GateKeeper and SHD. SneakySnake~\cite{alser2020sneakysnake} solves approximate string matching by converting it to a single net routing problem~\cite{lee1976use}. This enables SneakySnake to provide acceleration with and without hardware support. SneakySnake also improves the accuracy of SHD, GateKeeper, and Shouji. Lastly, GenASM~\cite{cali2020genasm} is a hardware-accelerated approximate string matching framework for genome sequence analysis with a modified Bitap algorithm~\cite{baeza1992new, wu1992fast}. When utilized for pre-alignment filtering of short reads, it provides a $3.7\times$ speedup over Shouji while improving the accuracy.
